\section{Prétraitement}

Toutes les images passent par le pipeline déterministe suivant avant toute extraction
de caractéristiques :

\begin{enumerate}
  \item \textbf{Conversion en niveaux de gris et mise à l'échelle.} Les images RGB
        sont converties en niveaux de gris ; les intensités sont ramenées à $[0, 1]$.
  \item \textbf{Recadrage des bordures noires.} Les zones sombres autour de la région
        cérébrale sont détectées et supprimées pour se concentrer sur le contenu anatomique.
  \item \textbf{Redimensionnement à \ImgSize{}.} Les images sont rééchantillonnées
        avec anti-crénelage, réduisant la dimensionnalité tout en préservant les textures fines.
  \item \textbf{Lissage gaussien} ($\sigma = 0{,}5$). Un filtre passe-bas atténue
        le bruit haute fréquence introduit lors de l'acquisition ou du redimensionnement.
  \item \textbf{Normalisation par image et CLAHE.} Chaque image est normalisée
        min--max individuellement, puis une égalisation adaptative de l'histogramme
        à contraste limité (CLAHE, clip limit $= 0{,}02$) améliore le contraste local.
\end{enumerate}

Cette séquence de prétraitement standardise la plage d'intensités entre patients
et appareils, et met en évidence les structures locales pertinentes pour les descripteurs
de texture et de gradient. La Figure~\ref{fig:preprocessed} montre un exemple de
coupe prétraitée.

\begin{figure}[H]
  \centering
  \includegraphics[width=0.45\linewidth,trim=130 520 130 70,clip]{descriptor_demo.png}
  \caption{Exemple de coupe IRM prétraitée (classe méningiome), \ImgSize{},
           après amélioration CLAHE.}
  \label{fig:preprocessed}
\end{figure}
